% -----------------------------------------------------------------------------
% Exact Ward--Takahashi identity from local U(1) covariance
% -----------------------------------------------------------------------------

局域 $U(1)$ 对称性不仅固定最小耦合顶角，也约束完整电子二点核与光子二点核。把精确一粒子不可约有效作用量记为 $\Gamma_{\rm 1PI}$，在经典背景四势 $A_\mu$ 中定义电子逆传播子
\begin{equation}
\boxed{
S_A^{-1}(x,y)
\equiv
\left.
\frac{\delta^2\Gamma_{\rm 1PI}}
{\delta\bar\psi(x)\,\delta\psi(y)}
\right|_{\psi=\bar\psi=0}.}
\label{eq:qed_background_inverse_propagator_definition}
\end{equation}
它是完整电子二点函数的逆核，并允许外部经典电磁背景任意强。采用
\[
A_\mu' = A_\mu-\partial_\mu\chi,
\qquad
U_\chi(x)=\exp\!\left(\frac{\ii \qe\chi(x)}{\hbar}\right),
\]
局域协变性要求
\begin{equation}
\boxed{
S_{A-\partial\chi}^{-1}(x,y)
=U_\chi(x)S_A^{-1}(x,y)U_\chi^{-1}(y).}
\label{eq:qed_inverse_propagator_gauge_covariance}
\end{equation}
这个关系不依赖微扰展开。规范固定只使光子二次核可逆，不改变费米子端点的局域相位变换律。

剥去电荷因子后，精确一粒子不可约电磁顶角定义为
\begin{equation}
\boxed{
\Gamma^\mu(x,y;z)
\equiv
-\frac{1}{\qe}
\left.
\frac{\delta S_A^{-1}(x,y)}{\delta A_\mu(z)}
\right|_{A=0}.}
\label{eq:qed_exact_vertex_functional_definition}
\end{equation}
树级时它退化为 $\Gamma^\mu=\gamma^\mu$。对\eqnrefs{eq:qed_inverse_propagator_gauge_covariance} 作无穷小规范变换，得到坐标空间 Ward--Takahashi 恒等式
\begin{equation}
\boxed{
\partial_\mu^z\Gamma^\mu(x,y;z)
=-\frac{\ii}{\hbar}
\bigl[\delta^{(4)}(z-x)-\delta^{(4)}(z-y)\bigr]
S^{-1}(x,y).}
\label{eq:qed_ward_takahashi_coordinate}
\end{equation}
两个接触项分别对应光子插入点与入射、出射费米子端点重合；相对负号来自两端相反的规范相位。

令入射电子动量为 $p^\mu$、出射电子动量为 $p^\mu+q^\mu$，按全文 $\exp(-\ii p\cdot x/\hbar)$ 的 Fourier 约定，\eqnrefs{eq:qed_ward_takahashi_coordinate} 等价于
\begin{equation}
\boxed{
q_\mu\Gamma^\mu(p+q,p)
=S^{-1}(p+q)-S^{-1}(p).}
\label{eq:qed_ward_takahashi_momentum}
\end{equation}
这里两边都是精确一粒子不可约核。任何固定圈阶都必须把 $S^{-1}$ 与 $\Gamma^\mu$ 按同一阶截断，才能保持这一约束。

在 $q\to0$ 且 $S^{-1}(p)$ 对外动量可微的区域，\eqnrefs{eq:qed_ward_takahashi_momentum} 给出
\begin{equation}
\boxed{
\Gamma^\mu(p,p)
=\frac{\partial S^{-1}(p)}{\partial p_\mu}.}
\label{eq:qed_ward_zero_momentum_vertex_derivative}
\end{equation}
因此电子波函数重整化与最小耦合顶角具有同一紫外乘法因子，
\begin{equation}
\boxed{Z_1=Z_2.}
\label{eq:qed_ward_z1_equals_z2}
\end{equation}
物理电荷归一化进一步固定在壳 Dirac 形状因子 $F_1(0)=1$；Pauli 形状因子 $F_2(0)$ 不由纵向 Ward 恒等式决定，必须由顶角的横向有限部分计算。

同一个局域对称性也约束光子二点的一粒子不可约量子修正。把真空极化张量定义为有效作用量中的二次核
\begin{equation}
\boxed{
\Gamma_{\Pi}^{(2)}[A]
=-\frac12\int\frac{\dd^4 q}{(2\pi\hbar)^4}\,
A_\mu(-q)\Pi^{\mu\nu}(q)A_\nu(q).}
\label{eq:qed-vacpol-quadratic-kernel-dp68}
\end{equation}
这里自由光子核和规范固定项不包含在 \(\Pi^{\mu\nu}\) 中。量子修正本身必须对 \(A_\mu\to A_\mu-\partial_\mu\chi\) 不变，因此
\begin{equation}
\boxed{q_\mu\Pi^{\mu\nu}(q)=0.}
\label{eq:qed_vacuum_polarization_transversality_from_1pi}
\end{equation}
在 Lorentz 不变真空中，横向性把它唯一约化为
\begin{equation}
\boxed{
\Pi^{\mu\nu}(q)
=\left(q^2\eta^{\mu\nu}-q^\mu q^\nu\right)\Pi(q^2).}
\label{eq:qed_vacuum_polarization_transverse_decomposition}
\end{equation}
因此真空极化的张量结构、$Z_1=Z_2$ 与外光子纵向极化解耦都是同一局域 $U(1)$ 约束的不同投影；具体圈积分只计算对称性未固定的标量函数和有限横向结构。
